%=======================================================================
%  ABBREVIATION CHART  (shared, appended to every chapter)
%=======================================================================
\needspace{6\baselineskip}
\T{Abbreviations used in these notes}

{\color{sub}\footnotesize Only these are ever shortened. Anything not on this list is written out in full.}

\vspace{2pt}
\begin{tabularx}{\textwidth}{@{}L{1.75cm}X@{\hspace{8pt}}L{1.75cm}X@{}}
\toprule
\hrow \thead{Short} & \thead{Means} & \thead{Short} & \thead{Means} \\
\midrule
BS      & base station                     & BTS     & base transceiver station \\
MS      & mobile station                   & MSC     & mobile switching centre \\
BSC     & base station controller           & MSK     & minimum shift keying \\
CCI     & co-channel interference           & ACI     & adjacent channel interference \\
SIR, S/I & signal to interference ratio     & SNR, S/N & signal to noise ratio \\
GoS     & grade of service                  & MAHO    & mobile assisted handoff \\
FCA     & fixed channel assignment          & DCA     & dynamic channel assignment \\
LOS     & line of sight                     & NLOS    & non line of sight \\
FDMA    & frequency division multiple access & TDMA   & time division multiple access \\
CDMA    & code division multiple access     & OFDM    & orthogonal frequency division multiplexing \\
FDD     & frequency division duplex         & TDD     & time division duplex \\
GSM     & global system for mobile communications & AMPS & advanced mobile phone system \\
PCS     & personal communication system     & DAS     & distributed antenna system \\
BCCH    & broadcast control channel         & EIRP    & effective isotropic radiated power \\
ERP     & effective radiated power          & RF      & radio frequency \\
Tx, TX  & transmitter                       & Rx, RX  & receiver \\
ISI     & intersymbol interference          & PDP     & power delay profile \\
DFE     & decision feedback equalizer       & MLSE    & maximum likelihood sequence estimation \\
LMS     & least mean square                 & RLS     & recursive least squares \\
MSE     & mean square error                 & MMSE    & minimum mean square error \\
ZF      & zero forcing                      & MRC     & maximal ratio combining \\
EGC     & equal gain combining              & BER     & bit error rate \\
PCM     & pulse code modulation             & DPCM    & differential PCM \\
ADPCM   & adaptive differential PCM         & LPC     & linear predictive coder \\
RPE-LTP & regular pulse excited, long term prediction & CELP & code excited linear prediction \\
RELP    & residual excited linear prediction & MPE    & multipulse excited \\
VSELP   & vector sum excited linear prediction & ATC   & adaptive transform coding \\
SBC     & sub-band coding                   & DCT     & discrete cosine transform \\
FEC     & forward error correction          & ARQ     & automatic repeat request \\
BCH     & Bose--Chaudhuri--Hocquenghem code & RS      & Reed--Solomon code \\
TCM     & trellis coded modulation          & PN      & pseudo-noise \\
attr(s) & attribute(s)                      & no.     & number \\
freq    & frequency                         & info    & information \\
ops     & operations                        & config  & configuration \\
rm      & remove                            & vs, Vs  & versus \\
w/, w/o & with, without                     & Eg      & example / for example \\
+Fig    & draw a figure                     & (I)     & importance / need / significance \\
$\rightarrow$ & leads to / therefore        & $+$     & and \\
\bottomrule
\end{tabularx}

\vspace{3pt}
{\color{sub}\footnotesize
\textbf{Exam-paper month codes} (used in the year tags):
Ba $=$ Baishakh $\cdot$ Bh $=$ Bhadra $\cdot$ Ash $=$ Ashwin $\cdot$
Ma $=$ Magh $\cdot$ Ch $=$ Chaitra. \\
\textbf{Bold} year tags are Regular exams, plain are Back exams.
\texttt{Monospace} year tags are BEI (EX 715); the rest are BEX (EX 751). \\
\textbf{Tier chips} over 22 papers, 2070--2082 BS:
\tS{} $=$ topic asked in 6 or more papers $\cdot$
\tF{} $=$ 3 to 5 papers $\cdot$
\tP{} $=$ 1 or 2 papers but worth $\ge$6 marks.
\textbf{Mark chips} such as \m{2+5} are the marks that question actually carried.
\pill{gdL}{gd}{syllabus, no PYQ yet} marks a syllabus topic that has never been examined but is
written up anyway. \added\ marks anything researched and added that was not in the source notes.}
